% =====================================================================
% TEMPLATE — MILP formulation (methodology asset A-P5-tpl)
% =====================================================================
% Purpose: present a MILP model (notation + model) in a consistent, readable
% style. Adapted from the "fleets" reference formulation, with ONE change:
% constraint blocks may be INTERLEAVED with their explanations, instead of
% listing every constraint first and all explanations at the end.
%
% Fill every <PLACEHOLDER>; delete guidance comments. Meant to be \input from a
% main.tex whose preamble provides:
%     \usepackage{amsmath,amssymb}
%     \usepackage{longtable}
%     \allowdisplaybreaks
%     \newcommand{\Sum}{\sum\limits_}   % optional shorthand
% (see the working harness in docs/<problem>/ — add \usepackage{longtable}
% there when this template is used).
%
% Conventions (keep them):
%   * NAMING (hard rule, for legibility — so a reader tells data from decisions
%     at a glance): PARAMETERS are ALWAYS UPPERCASE (multi-letter allowed, e.g.
%     $DUR^{MAX}_{dt}$, $V^{BOX}$); DECISION VARIABLES are ALWAYS lowercase
%     ($x_{bs}$, $y_b$). Sets are calligraphic ($\mathcal{X}$). Indices are
%     lowercase. When two parameters would collide on one letter, disambiguate
%     with sub/superscripts, never by lowercasing a parameter.
%   * Every constraint carries \label{eq:<name>} and its index range on the
%     right, separated with \hspace{10mm}.
%   * Big-M disjunctions read  ... - M (1 - <binary>).
%   * Objective and domain blocks stay as single blocks; the substantive
%     constraints are grouped and each group is explained right after it.
% =====================================================================

\part{MILP formulation}

% --- OPTIONAL scaffolding section -------------------------------------
% Include this section only if the model needs a prose preamble: the modeling
% scope, the discretization/abstraction chosen, and the key structural
% assumptions (each tagged, e.g. (A1), (A2)). Delete it for a bare notation +
% model write-up. When present it goes BEFORE \section{Notation}.
\section{Modeling scope and assumptions}
% 1–2 short paragraphs: what the model couples/decides and how the operation is
% discretized (e.g. into an ordered sequence, a set of positions/windows, ...);
% then a \begin{description} of the substantive assumptions.
<SCOPE: what the model couples/decides, and the chosen discretization/abstraction>
\begin{description}
  \item[\textbf{(A1) <assumption name>.}] <what it assumes and why it is safe>
  \item[\textbf{(A2) <assumption name>.}] <...>
\end{description}

% ---------------------------------------------------------------------
\section{Notation}\label{sec:<slug>-notation}
% ---------------------------------------------------------------------

\noindent \textit{Sets}

\begin{tabular}{p{22mm}lp{105mm}}
    $\mathcal{<SET1>}$                              & : & <description of the set> \\
    $\mathcal{<SET2>} \subset \mathcal{<SET1>}$     & : & <description; e.g. a subset indexed by an element> \\
    $\mathcal{<SET3>}_{<i>} \subset \mathcal{<SET2>}$ & : & <description of a parametrized family of sets> \\
\end{tabular}

\bigskip
\noindent \textit{Parameters}

\begin{longtable}{p{16mm}lp{108mm}}
    $<PARAM1>_{<i>}$        & : & <description, with units if any> \\
    $<PARAM2>_{<i><j>}$     & : & <description> \\
    $M$                     & : & Large enough constant for disjunctive (big-M) constraints. \\
\end{longtable}

\bigskip
\noindent \textit{Binary variables}

\begin{tabular}{p{10mm}lp{120mm}}
    $<x>_{<i>}$ & = & $\left\{ \begin{array}{ll}
                        1 & \mbox{if <condition holds for index $<i>$>} \\
                        0 & \mbox{otherwise}
                      \end{array} \right.$ \\[2mm]
    $<x2>_{<i><j>}$ & = & $\left\{ \begin{array}{ll}
                        1 & \mbox{if <condition>} \\
                        0 & \mbox{otherwise}
                      \end{array} \right.$ \\
\end{tabular}

\bigskip
\noindent \textit{Continuous / integer variables}

\begin{tabular}{p{14mm}lp{108mm}}
    $<y>_{<i>}$ & = & <meaning of the variable> \\
\end{tabular}

% ---------------------------------------------------------------------
\section{Model}\label{sec:<slug>-model}
% ---------------------------------------------------------------------

% --- Objective (single block) ----------------------------------------
\begin{align}
  \min \quad & <objective expression> \label{eq:<slug>-obj}
\end{align}
The objective~\eqref{eq:<slug>-obj} <one sentence: what it minimizes/maximizes
and what each term represents>.

\medskip
\noindent subject to:

% --- Constraint group 1: <name> --------------------------------------
\begin{align}
  & <lhs> = <rhs>
    \hspace{10mm} \forall\, <i> \in \mathcal{<SET1>}
    \label{eq:<slug>-<name1>}
\end{align}
Constraints~\eqref{eq:<slug>-<name1>} <what this group enforces and why>.

% --- Constraint group 2: <name> (big-M / multi-line example) ---------
\begin{align}
  & <lhs> \ge <rhs> - M\,\bigl(1 - <binary>\bigr) \notag \\
  & \hspace{10mm} \forall\, <i> \in \mathcal{<SET1>},\ <j> \in \mathcal{<SET2>}_{<i>}
    \label{eq:<slug>-<name2>}
\end{align}
Constraints~\eqref{eq:<slug>-<name2>} <explanation>. The big-M term deactivates
the bound unless $<binary> = 1$.

% --- Constraint group 3: <name> (range/counting example) -------------
\begin{align}
  & \sum_{<j> \in \mathcal{<SET2>}_{<i>}} <x2>_{<i><j>} \le <PARAM1>_{<i>}
    \hspace{10mm} \forall\, <i> \in \mathcal{<SET1>}
    \label{eq:<slug>-<name3>}
\end{align}
Constraints~\eqref{eq:<slug>-<name3>} <explanation>.

% --- Add further interleaved constraint groups here ------------------
% Pattern per group: a short \begin{align}...\end{align} with \label(s),
% immediately followed by a paragraph "Constraints~\eqref{...} ...".

% --- Variable domains (single block) ---------------------------------
\begin{align}
  & <x>_{<i>}     \in \{0,1\}            \hspace{10mm} \forall\, <i> \in \mathcal{<SET1>} \label{eq:<slug>-dom-x} \\
  & <x2>_{<i><j>} \in \{0,1\}            \hspace{10mm} \forall\, <i>,\, <j> \label{eq:<slug>-dom-x2} \\
  & <y>_{<i>}     \ge 0                  \hspace{10mm} \forall\, <i> \in \mathcal{<SET1>} \label{eq:<slug>-dom-y}
\end{align}
Constraints~\eqref{eq:<slug>-dom-x}--\eqref{eq:<slug>-dom-y} are the variable
domains.

% --- OPTIONAL scaffolding section -------------------------------------
% Include this section only if there are modeling caveats worth recording:
% linearity notes, big-M / horizon tightness, symmetry-breaking, or choices
% deferred to later phases. Delete it for a bare model. When present it goes
% AFTER the variable-domain block.
% ---------------------------------------------------------------------
\section{Notes and open modeling choices}\label{sec:<slug>-notes}
% ---------------------------------------------------------------------
\begin{itemize}
  \item \textbf{<note title>.} <e.g. why the model is fully linear / what sizes it well / big-M or horizon tightness>
  \item \textbf{<open choice>.} <a modeling choice deferred or left to tuning, and to which phase>
\end{itemize}
